\documentclass[conference]{IEEEtran}
\IEEEoverridecommandlockouts

% --- Standard IEEE Packages ---
\usepackage[utf8]{inputenc}
\usepackage[vietnamese]{babel}
\usepackage{amsmath,amssymb,amsfonts}
\usepackage{algorithmic}
\usepackage{graphicx}
\usepackage{textcomp}
\usepackage{xcolor}
\usepackage{booktabs}
\usepackage{multirow}
\usepackage{url}
\usepackage{cite}
\usepackage{hyperref}

\def\BibTeX{{\rm B\kern-.05em{\sc i\kern-.025em b}\kern-.08em
    T\kern-.1667em\lower.7ex\hbox{E}\kern-.125emX}}

\begin{document}

\title{Phân Loại Chọn Lọc và Hiệu Chuẩn Xác Suất Nhằm Nhận Diện Căng Thẳng Đáng Tin Cậy Từ Tín Hiệu Sinh Lý Đeo Cổ Tay}

\author{
    \IEEEauthorblockN{Phan Tuấn Đạt, Nguyễn Gia Huy, Trần Lợi Kiều Nhi, Nguyễn Minh Bảo}
    \IEEEauthorblockA{\textit{Khoa Công nghệ Thông tin / Trí tuệ Nhân tạo}, \textit{Trường Đại học FPT} \\
    Email: ptundatwork@gmail.com, huy522006@gmail.com, nhiklt2006@gmail.com, baonguyenminh17@gmail.com}
}

\maketitle

% =========================================================================
% TÓM TẮT (ABSTRACT)
% =========================================================================
\begin{abstract}
Các hệ thống theo dõi căng thẳng tâm lý liên tục qua thiết bị đeo cổ tay mở ra tiềm năng lớn trong chăm sóc sức khỏe chủ động. Tuy nhiên, cảm biến quang học thể tích sóng mạch (BVP) và dẫn truyền da (EDA) đặt tại cổ tay rất nhạy cảm với nhiễu động học. Những cử động tay tự nhiên, cử chỉ và dịch chuyển tư thế thường ngày kích hoạt các phản ứng giống như hệ thần kinh giao cảm (làm tăng nhịp tim và tiết mồ hôi do điều nhiệt), biểu hiện rất giống với phản ứng sinh lý khi căng thẳng tâm lý thực sự. Khi bị ép buộc phải đưa ra kết luận nhị phân trên 100\% các cửa sổ thời gian (Full Coverage), các mô hình học máy truyền thống tạo ra tỷ lệ báo động giả rất cao khi người dùng cử động, dẫn tới hiện tượng mệt mỏi vì chuông báo động (alarm fatigue) trong thực tế lâm sàng. Bài báo này đề xuất một khung nhận diện căng thẳng đáng tin cậy (Trustworthy AI), kết hợp thuật toán hiệu chuẩn xác suất hồi quy Isotonic với cơ chế phân loại chọn lọc (Selective Classification) tích hợp bối cảnh động học gia tốc trên bộ dữ liệu chuẩn WESAD ($N=15$). Cơ chế này cho phép mô hình từ chối đưa ra phán đoán (gán nhãn \textit{Bất định / Uncertain}) khi độ tin cậy hậu nghiệm không đạt ngưỡng an toàn $\tau$. Thực nghiệm đối đầu trên tám kiến trúc mô hình (bảy mô hình học máy dạng bảng và một mạng nơ-ron sâu MLP) dưới giao thức kiểm định nghiêm ngặt 15-fold Leave-One-Subject-Out (LOSO) kết hợp phân tầng nhóm đối tượng bên trong (GroupKFold) chứng minh rằng: lỗ hổng báo động giả khi vận động là vấn đề phổ biến trên mọi họ thuật toán ($\text{FAR}_{\text{motion}} \approx 12,2\% - 24,0\%$). Quy tắc từ chối đề xuất triệt tiêu từ 55,0\% đến 90,0\% các báo động giả này, đồng thời nâng cao điểm Macro $F_1$ thêm từ +3,6 đến +8,2 điểm phần trăm trên tập dữ liệu giữ lại. Tại ngưỡng tin cậy vận hành $\tau = 0,80$ với mô hình XGBoost, hệ thống đạt điểm Selective Macro $F_1$ gộp là \textbf{95,20\%} (trung bình theo đối tượng: $89,8\% \pm 17,6\%$) và điểm Stress $F_1$ gộp là \textbf{93,15\%} (trung bình theo đối tượng: $83,1\% \pm 32,6\%$, 95\% CI: [84,80\%, 98,54\%]), đưa sai số hiệu chuẩn ECE về mức rất thấp \textbf{0,0290} trên tập dữ liệu được giữ lại (giảm từ 0,0612 ở mô hình thô toàn phần), và giới hạn tỷ lệ báo động giả khi vận động $\text{FAR}_{\text{motion}}$ chỉ còn \textbf{4,07\%} (95\% CI: [0,0\%, 11,07\%]) trong khi vẫn duy trì độ bao phủ vận hành thực tế đạt $69,42\%$ (trung bình theo đối tượng: $69,6\% \pm 22,2\%$, 95\% CI: [58,34\%, 79,84\%]). Các phân tích chi tiết trên từng đối tượng, kiểm định phi tham số Wilcoxon, đối sánh dạng bảng - học sâu, và giải thích mô hình qua TreeSHAP đã làm sáng tỏ cơ chế sinh lý học và phân tích ca ngoại lai giảm phản ứng điện da (Subject S14).
\end{abstract}

\begin{IEEEkeywords}
Nhận diện Căng thẳng, Thiết bị Đeo Cổ tay, Bộ Dữ liệu WESAD, Phân loại Chọn lọc, Hiệu chuẩn Xác suất, Empatica E4, Báo động Giả do Vận động, Trí tuệ Nhân tạo Đáng tin cậy, Kiểm định LOSO.
\end{IEEEkeywords}

% =========================================================================
% I. GIỚI THIỆU
% =========================================================================
\section{Giới thiệu}
Căng thẳng tâm lý mãn tính (chronic stress) là một trong những yếu tố nguy cơ hàng đầu gây suy giảm miễn dịch, rối loạn chuyển hóa và các bệnh lý tim mạch nguy hiểm \cite{cohen2007psychological, schneiderman2005stress, can2019stress}. Sự phát triển mạnh mẽ của các thiết bị đeo thông minh thương mại (smartwatch, smartband) mở ra cơ hội theo dõi sinh lý liên tục, không xâm lấn trong đời sống hàng ngày \cite{gjoreski2016continuous, schmidt2018wesad, sano2018multimodal}.

Tuy nhiên, rào cản kỹ thuật lớn nhất khi triển khai các hệ thống này trong môi trường tự do (ambulatory monitoring) chính là \textit{hiện tượng báo động giả do vận động cơ học} (motion-induced false alarms) \cite{schumm2008effect, tamura2014wearable, zhang2014troika}. Những cử động tay tự nhiên, cử chỉ khi giao tiếp và dịch chuyển tư thế gây ra các biến dạng cơ học và dao động tự chủ rất giống với kích hoạt giao cảm \cite{boucsein2012electrodermal}. Dưới góc nhìn của cảm biến quang học thể tích sóng mạch (BVP) và dẫn truyền da (EDA) tại cổ tay, các đáp ứng sinh lý do vận động này biểu hiện gần như tương đồng với kích hoạt giao cảm do căng thẳng tâm lý cấp tính.

Trong các nghiên cứu trước đây trên bộ dữ liệu chuẩn WESAD \cite{schmidt2018wesad, siirtola2019continuous, bobade2020stress, dalmeida2021hrv, behinaein2021transformer, bottcher2023detection}, mô hình thường được đánh giá dưới cơ chế phán đoán bắt buộc 100\% (Full Coverage). Khi người đeo cử động tay hoặc thay đổi tư thế, các mô hình học máy truyền thống thường dự đoán sai với độ tự tin thái quá (overconfidence) \cite{guo2017calibration}. Hậu quả là tần suất báo động giả tăng vọt, gây ra hội chứng mệt mỏi vì chuông báo động (alarm fatigue) \cite{sendelbach2013alarm}, khiến người dùng mất niềm tin và tắt thiết bị theo dõi.

Để giải quyết thách thức này, nghiên cứu của chúng tôi tiếp cận bài toán theo nguyên lý \textbf{Học máy Đáng tin cậy (Trustworthy AI)}: thay vì ép buộc mô hình phải đoán bừa trên mọi cửa sổ dữ liệu, hệ thống cần lượng hóa độ bất định và từ chối phán đoán (Abstention) khi tín hiệu bị nhiễu hoặc nằm trong vùng mơ hồ thông qua phân loại chọn lọc (Selective Classification) \cite{chow1970optimum, cortes2016learning, geifman2017selective, geifman2019selectivenet, grandvalet2008support, hendrickx2024machine}.

\subsection{Các Câu hỏi Nghiên cứu (Research Questions - RQs)}
Nghiên cứu được định hướng giải quyết bốn câu hỏi khoa học trọng tâm:
\begin{itemize}
    \item \textbf{RQ1 (Tính nhạy cảm với chuyển động đa kiến trúc):} Hiện tượng báo động giả do vận động có tồn tại dai dẳng trên các họ thuật toán học máy khác nhau khi ép buộc phán đoán 100\% hay không?
    \item \textbf{RQ2 (Độ tin cậy và Hiệu chuẩn trong Phân loại Chọn lọc):} Cơ chế hiệu chuẩn xác suất kết hợp phân loại chọn lọc có thể triệt tiêu báo động giả do vận động đồng thời bảo toàn độ bao phủ vận hành cao hay không?
    \item \textbf{RQ3 (Học máy dạng bảng vs Học sâu):} Các mô hình cây (XGBoost, Random Forest, LightGBM) đối đầu với mạng nơ-ron sâu (MLP) như thế nào trên không gian đặc trưng sinh lý chuẩn hóa dưới giao thức LOSO?
    \item \textbf{RQ4 (Động lực học đặc trưng, Bối cảnh động học và Khái quát hóa ca ngoại lai):} Những tương quan đặc trưng sinh lý và động học nào kích hoạt cơ chế từ chối, và hệ thống phản ứng ra sao trước những cá nhân có phản ứng sinh lý không điển hình?
\end{itemize}

\subsection{Đóng góp Khoa học Chính}
Bài báo mang lại bốn đóng góp khoa học cụ thể:
\begin{enumerate}
    \item \textbf{Hình thức hóa và Định lượng Lỗ hổng Vận động ($\text{FAR}_{\text{motion}}$):} Chúng tôi định nghĩa chỉ số $\text{FAR}_{\text{motion}}$ để cô lập hành vi mô hình khi người dùng cử động và khảo sát toàn diện tám thuật toán (bảy mô hình bảng và một mạng MLP) dưới giao thức 15-fold Leave-One-Subject-Out (LOSO) \cite{saeb2017need, little2017using}. Kết quả chứng minh báo động giả khi vận động là lỗ hổng mang tính hệ thống ($\text{FAR}_{\text{motion}} \approx 12,2\% - 24,0\%$) trên mọi kiến trúc khi ép buộc phán đoán 100\%.
    \item \textbf{Quy trình Hiệu chuẩn Xác suất Hai cấp Chống Rò rỉ Dữ liệu:} Thiết kế đường ống hiệu chuẩn hồi quy Isotonic lồng ghép bên trong phân tầng nhóm đối tượng (\texttt{GroupKFold}). Quy trình này triệt tiêu hoàn toàn rò rỉ dữ liệu do cửa sổ trượt gối đầu 91,7\%, đưa sai số hiệu chuẩn ECE về mức \textbf{0,0290} trên tập dữ liệu giữ lại.
    \item \textbf{Khung Phân loại Chọn lọc Tích hợp Bối cảnh Động học Gia tốc:} Xây dựng quy tắc từ chối dựa trên ngưỡng tin cậy $\tau$ tích hợp 6 đặc trưng gia tốc kế (ACC). Khung giải pháp triệt tiêu từ 55,0\% đến 90,0\% báo động giả do vận động trên cả tám mô hình, giúp XGBoost đạt Selective Macro $F_1$ là \textbf{95,20\%} tại độ bao phủ $69,42\%$, được kiểm chứng qua 5 seed ngẫu nhiên và kiểm định Wilcoxon.
    \item \textbf{Giải thích Y sinh XAI và Phân tích Thực nghiệm Dạng bảng vs Học sâu:} Ứng dụng TreeSHAP \cite{lundberg2020local} lượng hóa đóng góp của 4 kênh cảm biến, làm rõ mẫu hình xung đột tín hiệu kích hoạt từ chối phán đoán, và phân tích sâu ca ngoại lai giảm phản ứng điện da (Subject S14).
\end{enumerate}

% =========================================================================
% II. TỔNG QUAN NGHIÊN CỨU & KHOẢNG TRỐNG KHOA HỌC
% =========================================================================
\section{Tổng quan nghiên cứu \& Khoảng trống khoa học}

\subsection{Nhận diện Căng thẳng Qua Thiết bị Đeo Cổ tay}
Các nghiên cứu ban đầu tập trung vào dây đeo ngực hoặc điện tâm đồ đa đạo trình (ECG) \cite{taskforce1996heart}. Với sự xuất hiện của các thiết bị đeo tay y tế đa cảm biến như Empatica E3 và E4 \cite{garbarino2014empatica}, các công trình gần đây chuyển dịch mạnh mẽ sang khai thác BVP, EDA, nhiệt độ da (TEMP) và gia tốc (ACC) \cite{can2019stress, gjoreski2016continuous}.

Schmidt cùng cộng sự \cite{schmidt2018wesad} công bố bộ dữ liệu chuẩn WESAD ($N=15$), chỉ ra rằng việc nhận diện stress từ cổ tay khó khăn hơn nhiều so với ngực do cổ tay cử động tự do, báo cáo độ chính xác nhị phân khoảng 87,8\% với mô hình Random Forest trên dữ liệu cổ tay (bao gồm cả gia tốc). Bobade và Vani \cite{bobade2020stress} khảo sát các mô hình truyền thống và học sâu CNN/LSTM; tuy nhiên quy trình đánh giá dùng K-Fold ngẫu nhiên gây rò rỉ dữ liệu nghiêm trọng giữa các cửa sổ trượt liền kề. Siirtola \cite{siirtola2019continuous} trích xuất biến thiên nhịp tim (HRV) từ BVP cổ tay, nhận định rằng sự dị biệt sinh lý giữa các cá nhân là rào cản lớn nhất. Can et al. \cite{can2019stress}, Sano et al. \cite{sano2018multimodal}, và Böttcher et al. \cite{bottcher2023detection} cùng xác nhận tính không đồng nhất sinh lý làm suy giảm mạnh độ chính xác. Gần đây, Behinaein et al. \cite{behinaein2021transformer} và Dalmeida và Masala \cite{dalmeida2021hrv} nghiên cứu các kiến trúc transformer trên ECG và HRV, tuy nhiên gánh nặng tính toán trên phần cứng nhúng vẫn là một thách thức lớn.

\subsection{Nhiễu Động học và Báo động Giả}
Chuyển động tay làm biến dạng sóng quang học BVP và làm xê dịch điện cực đo EDA \cite{tamura2014wearable, zhang2014troika, schumm2008effect}. Không chỉ gây biến dạng hình thái tín hiệu, vận động còn kích thích hệ giao cảm thật: nhịp tim tăng và tuyến mồ hôi tiết ra để điều hòa thân nhiệt \cite{boucsein2012electrodermal}.

Các nghiên cứu trước đây chủ yếu xử lý nhiễu bằng lọc số (Butterworth) hoặc triệt nhiễu thích nghi \cite{zhang2014troika}. Mặc dù lọc tần số loại được nhiễu điện cao tần, nó không thể phân tách được phản ứng căng thẳng tâm lý thật với các biến dạng tín hiệu do cử động tay gây ra. Khi mô hình bị ép buộc phán đoán 100\%, tỷ lệ báo động giả khi vận động luôn ở mức cao, dẫn tới hiện tượng mệt mỏi vì chuông báo \cite{sendelbach2013alarm}.

\subsection{Phân loại Chọn lọc và Học có Từ chối}
Bắt nguồn từ nền tảng của Chow \cite{chow1970optimum}, học có từ chối (learning with rejection) cho phép bộ phân loại bảo lưu ý kiến khi chi phí phán đoán sai lớn hơn chi phí từ chối. Geifman và El-Yaniv \cite{geifman2017selective, geifman2019selectivenet} mở rộng phân loại chọn lọc cho mạng nơ-ron sâu, phân tích đường cong rủi ro - độ bao phủ (Risk-Coverage). Các công trình lý thuyết của Cortes et al. \cite{cortes2016learning}, Grandvalet et al. \cite{grandvalet2008support}, và khảo sát của Hendrickx et al. \cite{hendrickx2024machine} đã khẳng định giá trị của cơ chế từ chối trong các bài toán nhạy cảm với rủi ro.

Mặc dù vậy, phân loại chọn lọc hầu như chưa được áp dụng bài bản cho bài toán nhận diện stress từ thiết bị đeo. Nghiên cứu này ứng dụng ngưỡng tin cậy $\tau$ kết hợp bối cảnh động học để chủ động triệt tiêu báo động giả do vận động.

\subsection{Hiệu chuẩn Xác suất và Rò rỉ Dữ liệu Cửa sổ Trượt}
Phân loại chọn lọc đòi hỏi xác suất hậu nghiệm phải được hiệu chuẩn chuẩn xác: độ tự tin dự đoán phải phản ánh đúng tần suất chính xác thực nghiệm \cite{guo2017calibration}. Các kỹ thuật hậu xử lý phổ biến gồm Platt scaling (Sigmoid) \cite{platt1999probabilistic}, Beta calibration \cite{kull2017beta}, và hồi quy Isotonic phi tham số \cite{zadrozny2002transforming, niculescu2005predicting}.

Trong chuỗi thời gian sinh lý, cửa sổ trượt 60 giây (bước nhảy 5 giây) tạo ra độ gối đầu tới 91,7\%. Nếu chia dữ liệu ngẫu nhiên, các cửa sổ liền kề của cùng một người sẽ lọt vào cả tập huấn luyện và tập hiệu chuẩn, gây ra ảo tưởng hiệu chuẩn (optimistic bias) \cite{saeb2017need, little2017using}. Chúng tôi giải quyết triệt để vấn đề này bằng thiết kế hai cấp lồng ghép phân tầng theo nhóm đối tượng.

\subsection{Lý do Ưu tiên Học máy Dạng bảng}
Việc ưu tiên các mô hình cây quyết định (Random Forest, LightGBM, XGBoost \cite{chen2016xgboost}) dựa trên ba căn cứ:
\begin{enumerate}
    \item \textbf{Quy mô Mẫu ($N=15$ Đối tượng):} Các công trình của Shwartz-Ziv và Armon \cite{shwartzziv2022tabular}, Grinsztajn et al. \cite{grinsztajn2022why}, và Borisov et al. \cite{borisov2024deep} chứng minh rằng trên dữ liệu dạng bảng quy mô vừa và nhỏ, mô hình cây thường vượt trội hoặc ngang ngửa học sâu. Mạng nơ-ron dung lượng lớn rất dễ overfit vào đặc điểm sinh lý cá nhân của số ít đối tượng, làm suy giảm khả năng tổng quát hóa LOSO.
    \item \textbf{Khả năng Giải thích Y sinh (XAI):} Các đặc trưng HRV và EDA (RMSSD, LF/HF, đỉnh SCR) gắn liền với sinh lý học thần kinh tự chủ \cite{taskforce1996heart, greco2016cvxeda, tarvainen2014kubios}. Khi kết hợp với TreeSHAP \cite{lundberg2020local}, mô hình cây cho phép bóc tách tường minh lý do đưa ra cảnh báo hoặc từ chối phán đoán.
    \item \textbf{Khả thi trên Phần cứng Nhúng Cổ tay:} Thiết bị đeo thương mại vận hành dưới giới hạn năng lượng vi điều khiển khắt khe. Mô hình cây kết hợp bảng tra cứu Isotonic tiêu tốn cực ít bộ nhớ và độ phức tạp suy luận $O(\text{depth})$, giúp triển khai trực tiếp trên chip mà không hao pin.
\end{enumerate}

Để đối chứng khách quan, chúng tôi huấn luyện thêm một mô hình mạng nơ-ron sâu MLP trên cùng không gian 25 đặc trưng chuẩn hóa dưới cùng điều kiện 15-fold LOSO. Bảng~\ref{tab:related_work_comparison} tóm tắt sự khác biệt giữa nghiên cứu này và các công trình liên quan.

\begin{table*}[t]
\caption{So sánh các nghiên cứu nhận diện căng thẳng qua thiết bị đeo và phương pháp đề xuất.}
\label{tab:related_work_comparison}
\centering
\resizebox{\textwidth}{!}{
\begin{tabular}{lcccccc}
\toprule
\textbf{Công trình} & \textbf{Thiết bị / Vị trí} & \textbf{Dòng Mô hình} & \textbf{Giao thức Đánh giá} & \textbf{Hiệu chuẩn Xác suất} & \textbf{Bối cảnh Động học (ACC)} & \textbf{Cơ chế Từ chối} \\
\midrule
Schmidt et al. (2018) \cite{schmidt2018wesad} & E4 (cổ tay) + RespiBAN (ngực) & LDA, Random Forest, AdaBoost & LOSO (15 folds) & Không & Đặc trưng cơ bản (~87,8\% Acc) & Không (Ép buộc 100\%) \\
Bobade \& Vani (2020) \cite{bobade2020stress} & E4 (cổ tay) + RespiBAN & SVM, KNN, CNN, LSTM & K-Fold Ngẫu nhiên & Không & Đặc trưng cơ bản & Không (Ép buộc 100\%) \\
Siirtola (2019) \cite{siirtola2019continuous} & Smartwatch (cổ tay - BVP) & Random Forest, QDA & LOSO & Không & Không đưa vào & Không (Ép buộc 100\%) \\
Gjoreski et al. (2016) \cite{gjoreski2016continuous} & Thiết bị cổ tay & Random Forest, Ensemble & Phòng lab \& Đời thực & Không & Lọc hoạt động thô & Không (Ép buộc 100\%) \\
Can et al. (2019) \cite{can2019stress} & Thiết bị đeo đa cảm biến & SVM, Gradient Boosting & Đa giao thức & Không & Phân tích khảo sát & Không (Ép buộc 100\%) \\
Behinaein et al. (2021) \cite{behinaein2021transformer} & Thiết bị đeo ECG & Transformer, Học sâu & LOSO & Không & Không đưa vào & Không (Ép buộc 100\%) \\
\midrule
\textbf{Nghiên cứu Đề xuất} & \textbf{Empatica E4 (cổ tay)} & \textbf{8 Mô hình (7 Dạng bảng + MLP)} & \textbf{LOSO Hai cấp 15-folds} & \textbf{Có (Inner GroupKFold Isotonic)} & \textbf{Tích hợp 6 Đặc trưng ACC} & \textbf{Có (Quy tắc Chọn lọc $\tau = 0,80$)} \\
\bottomrule
\end{tabular}}
\end{table*}

% =========================================================================
% III. PHƯƠNG PHÁP ĐỀ XUẤT
% =========================================================================
\section{Phương pháp đề xuất}

\subsection{Quy trình Tiền xử lý \& Trích xuất Đặc trưng}
Dữ liệu liên tục từ vòng đeo tay Empatica E4 được phân đoạn thành các cửa sổ trượt $W = 60$\,s với bước nhảy $\Delta t = 5$\,s (độ gối đầu 91,7\%). Mặc dù khuyến nghị Task Force tiêu chuẩn yêu cầu 5 phút cho phân tích HRV thông thường \cite{taskforce1996heart}, cửa sổ 60 giây thuộc nhóm HRV siêu ngắn hạn (Ultra-Short-Term HRV - UST-HRV) đã được chứng minh là khả thi trong giám sát di động \cite{shaffer2017overview}:
\begin{itemize}
    \item \textbf{Tín hiệu BVP (64\,Hz):} Lọc dải thông Butterworth bậc 2 (0,5--8,0\,Hz). Dò tìm các đỉnh tâm thu để xác định chuỗi khoảng cách giữa hai nhịp tim liên tiếp (IBI). Loại bỏ các khoảng IBI nhiễu ngoài dải 300--1800\,ms, sau đó nội suy spline bậc ba về tần số 4\,Hz \cite{tarvainen2014kubios}. Các đặc trưng trích xuất gồm miền thời gian (Mean HR, SDNN, RMSSD, pNN50) và miền tần số qua phương pháp Welch: công suất tần số thấp LF (0,04--0,15\,Hz), tần số cao HF (0,15--0,40\,Hz), và tỷ số cân bằng giao cảm/phó giao cảm LF/HF.
    \item \textbf{Tín hiệu EDA (4\,Hz):} Phân tách bằng thuật toán cvxEDA \cite{greco2016cvxeda} thành thành phần trương lực (tonic SCL) và thành phần pha (phasic SCR). Trích xuất số lượng đỉnh SCR, biên độ đỉnh cực đại, giá trị trung bình và phương sai của cả tonic và phasic, cùng năng lượng pha $\sum \text{Phasic}^2$ \cite{boucsein2012electrodermal}.
    \item \textbf{Tín hiệu TEMP (4\,Hz):} Nhiệt độ bề mặt da phản ánh sự co mạch ngoại vi. Trích xuất giá trị trung bình, độ lệch chuẩn, các cực trị min/max, và hệ số góc hồi quy tuyến tính.
    \item \textbf{Tín hiệu ACC (32\,Hz):} Gia tốc ba trục được tính độ lớn vector Euclidean $a(t) = \sqrt{a_x^2(t) + a_y^2(t) + a_z^2(t)}$. Trích xuất đủ 6 đặc trưng: giá trị trung bình, độ lệch chuẩn, phương sai $\text{Var}(\text{ACC})$ (đại diện cho cường độ rung lắc cơ học), giá trị đỉnh cực đại, tỷ lệ qua điểm không (ZCR), và năng lượng phổ tín hiệu.
\end{itemize}

Tập dữ liệu hoàn chỉnh gồm 5.294 cửa sổ quan sát trên toàn bộ 15 đối tượng hợp lệ của WESAD ($S2, S3, \dots, S17$; trong đó $S1$ là phiên chạy thử nghiệm giao thức và $S12$ bị lỗi phần cứng cảm biến đã được nhóm tác giả WESAD loại bỏ từ đầu \cite{schmidt2018wesad}). Không gian đặc trưng chuẩn hóa gồm đúng 25 đặc trưng liên tục.

\subsection{Hiệu chuẩn Xác suất Hậu nghiệm}
Gọi $x \in \mathbb{R}^d$ là vector đặc trưng và $y \in \{0, 1\}$ là nhãn thực tế (0: Không Stress, 1: Căng thẳng). Đầu ra thô $f(x)$ của mô hình được chuyển đổi thành xác suất hiệu chuẩn qua hàm hồi quy Isotonic đơn điệu \cite{zadrozny2002transforming, niculescu2005predicting}:
\begin{equation}
\hat{P}(Y=1|x) = g(f(x))
\end{equation}
trong đó $g(\cdot)$ tối thiểu hóa hàm mất mát bình phương dưới ràng buộc đơn điệu:
\begin{equation}
\begin{aligned}
\min_{g} \ & \sum_{i=1}^{N_{\text{cal}}} \left( y_i - g(f(x_i)) \right)^2 \\
\text{s.t.} \ & g(f(x_i)) \le g(f(x_j)) \quad \forall \ f(x_i) \le f(x_j)
\end{aligned}
\end{equation}

\subsection{Quy tắc Phán đoán Chọn lọc}
Gọi $\hat{p}_1 = \hat{P}(Y=1|x)$. Chúng tôi thiết lập quy tắc phán đoán an toàn dựa trên ngưỡng tin cậy $\tau \in [0,5; 1,0)$:
\begin{equation}
\hat{y}(x) =
\begin{cases}
1 \ (\text{Căng thẳng}), & \text{nếu } \hat{p}_1 \ge \tau \\
0 \ (\text{Bình thường}), & \text{nếu } \hat{p}_1 \le 1 - \tau \\
\text{Bất định}, & \text{nếu } 1 - \tau < \hat{p}_1 < \tau
\end{cases}
\end{equation}
Khi xác suất hậu nghiệm rơi vào khoảng lưỡng lự $(1-\tau, \tau)$, hệ thống bảo lưu quyết định (Abstain). Cơ chế này chủ động gạt bỏ các mẫu nằm ở vùng biên phán đoán vốn là nguồn cơn gây ra phần lớn các báo động giả khi người dùng cử động.

Chúng tôi đối chuẩn quy tắc chọn lọc này với hai phương pháp cơ sở: (i) cổng cơ học đơn giản ($\text{Từ chối nếu } \text{Var}(\text{ACC}) \ge \theta_{\text{motion}}$), và (ii) từ chối ngẫu nhiên tại cùng mức bao phủ ($69,42\%$). Mặc dù cổng cơ học triệt tiêu hoàn toàn báo động giả ($\text{FAR}_{\text{motion}} = 0,0\%$), nó loại bỏ một cách mù quáng tới 40,6\% các cửa sổ căng thẳng thực sự (764 trên tổng số 1.882 cửa sổ stress) trong các pha thuyết trình và tính nhẩm TSST, khiến Stress $F_1$ sụt giảm nghiêm trọng xuống 84,16\% (Macro $F_1 = 89,36\%$). Từ chối ngẫu nhiên đạt Stress $F_1 = 86,69\%$ và Macro $F_1 = 89,80\%$ ($\text{FAR}_{\text{motion}} = 6,78\%$). Ngược lại, quy tắc chọn lọc đề xuất bảo toàn trọn vẹn các giai đoạn stress thật trong khi cử động tay, đạt Stress $F_1 = 93,15\%$ và Macro $F_1 = 95,20\%$, đồng thời ghìm $\text{FAR}_{\text{motion}}$ xuống $4,07\%$.

% =========================================================================
% IV. THIẾT LẬP THỰC NGHIỆM
% =========================================================================
\section{Thiết lập thực nghiệm}

\subsection{Giao thức Kiểm định Nghiêm ngặt LOSO Hai cấp}
Để đánh giá khả năng tổng quát hóa trên người chưa từng xuất hiện trong tập huấn luyện, chúng tôi triển khai giao thức 15-fold Leave-One-Subject-Out (LOSO). Ở mỗi fold $k \in \{1, \dots, 15\}$, toàn bộ dữ liệu của một đối tượng $S_k$ được tách riêng làm tập kiểm thử $\mathcal{D}_{\text{test}}^{(k)}$.

Để hiệu chuẩn xác suất mà không làm rò rỉ dữ liệu, 14 đối tượng còn lại $\mathcal{D}_{\text{train}}^{(k)}$ được chia tiếp bằng kiểm định chéo nội bộ 3-fold \texttt{GroupKFold} phân tầng chặt chẽ theo \texttt{subject\_id}. Các dự đoán ngoài tập (out-of-fold) của vòng trong được dùng để huấn luyện bộ hiệu chuẩn Isotonic $g^{(k)}$. Sau đó, toàn bộ đường ống đã hiệu chuẩn mới được áp dụng lên đối tượng kiểm thử độc lập $\mathcal{D}_{\text{test}}^{(k)}$.

\subsection{Tập Mô hình Khảo sát (7 Dạng bảng + 1 Học sâu MLP)}
Tám thuật toán đại diện cho các trường phái học máy khác nhau được đưa vào khảo sát:
\begin{enumerate}
    \item \textbf{Logistic Regression (LR):} Mô hình tuyến tính với điều chuẩn L2 ($\text{max\_iter}=500$).
    \item \textbf{Decision Tree (CART):} Cây quyết định đơn lẻ với $\text{max\_depth}=5$.
    \item \textbf{k-Nearest Neighbors (kNN):} Phân loại láng giềng gần nhất ($k=5$).
    \item \textbf{AdaBoost:} Tăng cường tương thích với 50 gốc cây quyết định.
    \item \textbf{Random Forest (RF):} Rừng ngẫu nhiên gồm 100 cây ($\text{max\_depth}=5$).
    \item \textbf{LightGBM:} Gradient boosting (100 cây, $\text{max\_depth}=5$, $\text{learning\_rate}=0,08$).
    \item \textbf{MLP (Học sâu):} Mạng Perceptron đa tầng (kiến trúc: 25--64--32--2, 3.810 tham số, kích hoạt ReLU, tối ưu Adam, tốc độ học $10^{-3}$, suy giảm trọng số $10^{-4}$, dừng sớm dựa trên tập validation nội bộ để chống rò rỉ).
    \item \textbf{XGBoost (Đề xuất):} Gradient boosting (100 cây, $\text{max\_depth}=5$, tốc độ học $0,08$, lấy mẫu dòng $0,8$, lấy mẫu cột $0,8$).
\end{enumerate}

\subsection{Chỉ số Đánh giá và Khoảng Tin cậy Bootstrap}
Hiệu năng được đo lường bằng Macro $F_1$, $F_1$ của lớp Căng thẳng (Stress), và Độ bao phủ hữu dụng (Usable Coverage):
\begin{equation}
\text{Coverage} = \frac{|\{x \in \mathcal{D}_{\text{test}} : \hat{y}(x) \neq \text{Bất định}\}|}{|\mathcal{D}_{\text{test}}|}
\end{equation}
Sai số hiệu chuẩn kỳ vọng ECE được tính trên $M=10$ khoảng xác suất bằng nhau:
\begin{equation}
\text{ECE} = \sum_{m=1}^M \frac{|B_m|}{N} \left| \text{acc}(B_m) - \text{conf}(B_m) \right|
\end{equation}

So sánh thống kê được thực hiện bằng kiểm định cặp Wilcoxon signed-rank kèm hiệu chỉnh Holm-Bonferroni ($p_{\text{holm}}$). Khoảng tin cậy 95\% được ước tính qua phương pháp Bootstrap theo cụm đối tượng với $B=1.000$ lần lặp lại.

\subsection{Chỉ số Báo động Giả khi Vận động ($\text{FAR}_{\text{motion}}$)}
Để đo lường tỷ lệ báo động giả khi vận động mà không làm tràn cột hiển thị, gọi $\mathcal{M}_{\text{test}} = \{x \in \mathcal{D}_{\text{test}} \mid y=0, \ \text{Var}(\text{ACC}) \ge \theta_{\text{motion}}\}$ là tập hợp các cửa sổ bình thường (không căng thẳng) nhưng có vận động tay mạnh. Chỉ số $\text{FAR}_{\text{motion}}$ được định nghĩa như sau:
\begin{equation}
\text{FAR}_{\text{motion}} = \frac{|\{x \in \mathcal{M}_{\text{test}} : \hat{y}(x) = 1\}|}{|\{x \in \mathcal{M}_{\text{test}} : \hat{y}(x) \neq \text{Bất định}\}|}
\end{equation}
với $\theta_{\text{motion}} = 15,385$\,$\text{m}^2/\text{s}^4$ là phân vị thứ 80 của phương sai gia tốc trên toàn bộ tập dữ liệu WESAD (1.059/5.294 cửa sổ), cô lập nhóm 20% các cửa sổ có cử động tay rõ rệt ($N=295$ cửa sổ vận động không stress thuộc tập $\mathcal{M}_{\text{test}}$).

% =========================================================================
% V. KẾT QUẢ THỰC NGHIỆM VÀ THẢO LUẬN
% =========================================================================
\section{Kết quả thực nghiệm và Thảo luận}

\subsection{Khảo sát Đa Mô hình Benchmark (Trả lời RQ1)}

Bảng~\ref{tab:benchmark_8_models} tóm tắt hiệu năng của cả tám mô hình dưới hai chế độ: bắt buộc phán đoán 100\% và phân loại chọn lọc ($\tau = 0,80$).

\begin{table*}[t]
\caption{Hiệu năng đối đầu giữa tám kiến trúc mô hình (bảy mô hình dạng bảng và một mạng sâu MLP) dưới kiểm định 15-fold Leave-One-Subject-Out (LOSO) trên bộ dữ liệu WESAD.}
\label{tab:benchmark_8_models}
\centering
\resizebox{\textwidth}{!}{
\begin{tabular}{lcccc|cc|ccccc}
\toprule
\multirow{2}{*}{\textbf{Kiến trúc Mô hình}} & \multicolumn{4}{c|}{\textbf{Bắt buộc Toàn phần (100\%) Thô}} & \multicolumn{2}{c|}{\textbf{Đã Hiệu chuẩn Isotonic}} & \multicolumn{5}{c}{\textbf{Phân loại Chọn lọc ($\tau = 0,80$)}} \\
& \textbf{Độ CX (\%)} & \textbf{Macro $F_1$} & \textbf{ECE} & $\text{FAR}_{\text{mot}}$ & \textbf{ECE} & $\text{FAR}_{\text{mot}}$ & \textbf{Bao phủ} & \textbf{Độ CX (\%)} & \textbf{Macro $F_1$} & \textbf{ECE} & $\text{FAR}_{\text{mot}}$ & \textbf{Giảm FAR} \\
\midrule
Logistic Regression & 88,93 & 88,00\% & 0,0460 & 23,05\% & 0,0469 & 24,41\% & 72,48\% & 93,54 & 92,23\% & 0,0485 & 10,17\% & 55,9\% \\
Decision Tree (CART)& 82,43 & 81,09\% & 0,1468 & 22,71\% & 0,0973 & 11,86\% & 50,70\% & 91,73 & 85,94\% & 0,0465 & 3,05\% & 86,6\% \\
k-Nearest Neighbors & 86,95 & 85,48\% & 0,0835 & 20,34\% & 0,0548 & 20,34\% & 64,94\% & 95,03 & 93,69\% & 0,0500 & 9,15\% & 55,0\% \\
AdaBoost Classifier & 90,52 & 89,66\% & 0,2727 & 13,90\% & 0,0722 & 12,88\% & 77,13\% & 94,24 & 93,29\% & 0,0548 & 1,69\% & 87,8\% \\
Random Forest       & 89,21 & 88,01\% & 0,0638 & 13,90\% & 0,0390 & 16,61\% & 75,03\% & 95,32 & 94,49\% & 0,0336 & 3,39\% & 75,6\% \\
LightGBM            & 86,53 & 85,66\% & 0,0915 & 13,56\% & 0,0686 & 9,49\%  & 63,15\% & 95,09 & 93,65\% & 0,0333 & 1,36\% & \textbf{90,0\%} \\
MLP (Học sâu)       & 89,63 & 88,95\% & 0,0312 & 24,00\% & 0,0336 & 25,63\% & 75,16\% & 95,76 & 95,38\% & 0,0179 & 10,17\% & 57,6\% \\
\textbf{XGBoost (Đề xuất)} & \textbf{88,97} & \textbf{88,13\%} & \textbf{0,0612} & \textbf{12,20\%} & \textbf{0,0744} & \textbf{10,85\%} & \textbf{69,42\%} & \textbf{96,08} & \textbf{95,20\%} & \textbf{0,0290} & \textbf{4,07\%} & \textbf{66,6\%} \\
\bottomrule
\end{tabular}}
\end{table*}

\begin{figure*}[t]
\centering
\includegraphics[width=0.98\linewidth]{Images/fig_multimodel_benchmark.png}
\caption{So sánh hiệu năng đa mô hình trên tám thuật toán: (A) Tỷ lệ báo động giả khi vận động ($\text{FAR}_{\text{motion}}$) được dập tắt từ 55,0\% đến 90,0\% trên mọi mô hình; (B) Mức tăng trưởng điểm Selective Macro $F_1$ tại ngưỡng $\tau = 0,80$.}
\label{fig:multimodel_benchmark}
\end{figure*}

Kết quả benchmark chỉ ra hai phát hiện cốt lõi:
\begin{enumerate}
    \item \textbf{Lỗ hổng Báo động Giả Phổ biến:} Dưới chế độ bắt buộc phán đoán 100\%, tất cả các mô hình đều ghi nhận tỷ lệ báo động giả khi vận động rất cao: $\text{FAR}_{\text{motion}}$ dao động từ $12,20\%$ (XGBoost) đến $24,00\%$ (MLP). Mặc dù MLP đạt điểm tổng thể cạnh tranh ($F_1 = 88,95\%$), mô hình này lại có tỷ lệ báo động giả khi vận động cao nhất ($24,00\%$). Điều này khẳng định hiện tượng báo động giả là do nhiễu loạn sinh lý khách quan chứ không phải nhược điểm riêng của một thuật toán cụ thể.
    \item \textbf{Hiệu quả của Phân loại Chọn lọc:} Áp dụng cơ chế chọn lọc tại $\tau = 0,80$ giúp dập tắt mạnh mẽ báo động giả trên \textbf{toàn bộ tám kiến trúc} (mức giảm tương đối từ 55,0\% đến 90,0\%), đồng thời nâng Macro $F_1$ thêm từ +3,6 đến +8,2 điểm phần trăm trên tập dữ liệu giữ lại (ví dụ kNN tăng +8,21 điểm phần trăm từ 85,48\% lên 93,69\%). XGBoost đạt điểm cân bằng Pareto tối ưu: giới hạn $\text{FAR}_{\text{motion}}$ xuống còn $4,07\%$ (giảm $66,6\%$ so với baseline thô ban đầu) và đạt Selective Macro $F_1$ là 95,20\% tại độ bao phủ $69,42\%$. Mặc dù LightGBM đạt $\text{FAR}_{\text{motion}}$ thấp nhất (1,36\%, giảm 90,0\%), mô hình này phải trả giá bằng độ bao phủ thấp hơn nhiều (63,15\% so với 69,42\%) và Macro $F_1$ thấp hơn (93,65\% so với 95,20\%). XGBoost được ưu tiên vì giữ lại nhiều hơn 6,27 điểm phần trăm cửa sổ theo dõi hữu dụng (69,42\% so với 63,15\%).
\end{enumerate}

\subsection{Học máy Dạng bảng (XGBoost) vs Học sâu (MLP) (Trả lời RQ3)}
Kiểm định Wilcoxon signed-rank và khoảng tin cậy bootstrap trên 15 fold LOSO cho thấy:
\begin{itemize}
    \item \textbf{So sánh khi Bắt buộc Phán đoán 100\%:} Macro $F_1$ giữa MLP ($88,95\%$) và XGBoost ($88,13\%$) không có sự khác biệt mang ý nghĩa thống kê ($W = 52,0, n_{\text{nonzero}} = 14, p = 0,975, r_{rb} = -0,010, p_{\text{holm}} = 1,0$). Tuy nhiên, XGBoost có độ nhạy cảm với vận động thấp hơn đáng kể ($\text{FAR}_{\text{motion}} = 12,20\%$ so với $24,00\%$ của MLP).
    \item \textbf{Hiệu năng Chọn lọc tại $\tau = 0,80$:} Cả hai mô hình đều vượt mốc 95\% Selective Macro $F_1$ ($95,38\%$ với MLP và $95,20\%$ với XGBoost). Khoảng tin cậy bootstrap cặp theo cụm đối tượng $(MLP - XGB)$ đều chứa số không: $\Delta F_1 = +0,0117$ (95\% CI: $[-0,0256; +0,0513]$) và $\Delta \text{FAR}_{\text{sel}} = +2,03\%$ (95\% CI: $[-5,44\%; +8,54\%]$), chứng minh cơ chế chọn lọc vận hành ổn định trên cả hai trường phái.
    \item \textbf{Tính Ổn định Giữa các Đối tượng:} Trên cá nhân S14 (Fold 12), MLP bắt được các tương quan phi tuyến tinh tế ($F_{1,\text{stress}} = 97,2\%$), trong khi các mô hình cây bị suy giảm do biên độ EDA bị triệt tiêu. Ngược lại, MLP lại từ chối phán đoán hoàn toàn trên đối tượng S17 (Fold 15), dẫn đến độ bao phủ 0\%. Trong khi đó, XGBoost duy trì độ bao phủ ổn định trên toàn bộ 15 đối tượng (gộp $69,42\%$, trung bình $69,6\% \pm 22,2\%$).
    \item \textbf{Khía cạnh Phần cứng Thiết bị Đeo:} Dưới góc độ phần cứng nhúng, cả XGBoost và mạng MLP nhỏ (khoảng 3.800 tham số) đều có thời gian suy luận cực nhanh ($< 1$\,ms) trên vi điều khiển ARM Cortex-M4/M33. Tuy nhiên, mô hình cây có lợi thế diễn giải trực tiếp qua các ngưỡng phân nhánh. Trên thực tế, nút thắt tính toán chính trên thiết bị đeo nằm ở khâu trích xuất đặc trưng sinh lý (như tối ưu lồi cvxEDA), vốn có thể được lập lịch định kỳ để tiết kiệm pin.
\end{itemize}

\subsection{Sự Đánh đổi Chọn lọc và Đánh giá Chi tiết 15 Folds (Trả lời RQ2 \& RQ3)}

Bảng~\ref{tab:ketqua_full} thể hiện sự đánh đổi có hệ thống của XGBoost qua các ngưỡng $\tau \in [0,60; 0,90]$. Mức giảm báo động giả được minh họa ở Hình~\ref{fig:far_motion}, giản đồ tin cậy hiệu chuẩn xác suất ở Hình~\ref{fig:calibration}, và đường cong rủi ro - độ bao phủ ở Hình~\ref{fig:risk_coverage}.

\begin{table}[htbp]
\caption{Đặc tính đánh đổi trong phân loại chọn lọc của XGBoost qua các ngưỡng tin cậy $\tau$ dưới kiểm định 15-fold LOSO.}
\label{tab:ketqua_full}
\centering
\resizebox{\linewidth}{!}{
\begin{tabular}{lccccc}
\toprule
\textbf{Cấu hình} & \textbf{Độ Bao phủ} & \textbf{Macro $F_1$} & \textbf{Stress $F_1$} & \textbf{ECE} & $\text{FAR}_{\text{motion}}$ \\
\midrule
XGBoost (Thô Toàn phần 100\%)     & 100,0\% & 88,13\% & 84,99\% & 0,0612 & 12,20\% \\
XGBoost + Isotonic (100\%)       & 100,0\% & 89,88\% & 86,68\% & 0,0744 & 10,85\% \\
\midrule
Chọn lọc ($\tau=0,60$)          & 89,76\% & 92,86\% & 90,45\% & 0,0534 & 7,46\% \\
Chọn lọc ($\tau=0,70$)          & 81,79\% & 93,68\% & 91,28\% & 0,0406 & 4,41\% \\
\textbf{Chọn lọc Vận hành ($\tau=0,80$)} & \textbf{69,42\%} & \textbf{95,20\%} & \textbf{93,15\%} & \textbf{0,0290} & \textbf{4,07\%} \\
Chọn lọc Khắt khe ($\tau=0,90$) & 50,38\% & 96,93\% & 95,54\% & 0,0161 & 0,68\% \\
\bottomrule
\end{tabular}}
\end{table}

\begin{figure}[htbp]
\centering
\includegraphics[width=0.98\linewidth]{Images/far_motion_reduction.png}
\caption{Mức độ dập tắt tỷ lệ báo động giả khi vận động ($\text{FAR}_{\text{motion}}$) từ chế độ bắt buộc sang phân loại chọn lọc.}
\label{fig:far_motion}
\end{figure}

\begin{figure}[htbp]
\centering
\includegraphics[width=0.98\linewidth]{Images/calibration_curve.png}
\caption{Giản đồ độ tin cậy: Hồi quy Isotonic đưa phân phối xác suất dự đoán của XGBoost tiệm cận đường tham chiếu lý tưởng.}
\label{fig:calibration}
\end{figure}

\begin{figure}[htbp]
\centering
\includegraphics[width=0.98\linewidth]{Images/risk_coverage_curve.png}
\caption{Đường cong đánh đổi rủi ro - độ bao phủ (Risk-Coverage) qua các ngưỡng $\tau \in [0,50; 0,95]$. Rủi ro phân loại chọn lọc giảm đơn điệu khi ngưỡng tin cậy $\tau$ tăng.}
\label{fig:risk_coverage}
\end{figure}

Bảng~\ref{tab:per_subject_loso} và Hình~\ref{fig:subject_loso_fold} trình bày chi tiết kết quả trên từng đối tượng trong số 15 người tham gia WESAD ($S2 \rightarrow S17$, không tính đối tượng thử nghiệm $S1$ và đối tượng lỗi cảm biến $S12$).

\begin{table*}[t]
\caption{Kết quả chi tiết trên 15 fold Leave-One-Subject-Out (LOSO) với mô hình XGBoost hiệu chuẩn chọn lọc ($\tau = 0,80$), gắn nhãn theo mã định danh gốc của WESAD (Subject S14 đánh giá ở Fold 12; Subject S12 bị loại bỏ do lỗi cảm biến WESAD).}
\label{tab:per_subject_loso}
\centering
\resizebox{\textwidth}{!}{
\begin{tabular}{lcccc|ccc|cccc}
\toprule
\multirow{2}{*}{\textbf{Fold / Đối tượng}} & \multirow{2}{*}{\textbf{Cửa sổ $n$}} & \textbf{Tỷ lệ Stress} & \textbf{Tỷ lệ Vận động} & \multicolumn{3}{c|}{\textbf{Phán đoán Bắt buộc (Full Coverage)}} & \multicolumn{4}{c}{\textbf{Phân loại Chọn lọc ($\tau = 0,80$)}} \\
& & \textbf{(\%)} & \textbf{(\%)} & \textbf{$F_{1,\text{XGB}}$ (\%)} & \textbf{$F_{1,\text{Cal}}$ (\%)} & \textbf{$\text{FAR}_{\text{Cal}}$ (\%)} & \textbf{Độ Bao phủ (\%)} & \textbf{$F_{1,\text{sel}}$ (\%)} & \textbf{Độ CX$_{\text{sel}}$ (\%)} & \textbf{$\text{FAR}_{\text{sel}}$ (\%)} \\
\midrule
Fold 1 (S2)  & 338 & 34,3\% & 18,9\% & 76,6\% & 71,8\% & 0,0\% & 73,1\% & 17,6\% & 88,7\% & 0,0\% \\
Fold 2 (S3)  & 341 & 35,5\% & 24,0\% & 73,1\% & 71,2\% & 63,6\% & 59,8\% & 81,1\% & 77,0\% & 54,5\% \\
Fold 3 (S4)  & 345 & 34,8\% & 15,1\% & 100,0\% & 100,0\% & 0,0\% & 100,0\% & 100,0\% & 100,0\% & 0,0\% \\
Fold 4 (S5)  & 353 & 34,3\% & 5,4\% & 89,0\% & 92,1\% & 78,9\% & 71,7\% & 96,0\% & 97,6\% & 31,6\% \\
Fold 5 (S6)  & 351 & 35,0\% & 4,6\% & 81,7\% & 94,2\% & 7,1\% & 53,8\% & 100,0\% & 100,0\% & 0,0\% \\
Fold 6 (S7)  & 351 & 34,5\% & 17,4\% & 85,5\% & 90,6\% & -- & 60,4\% & 98,6\% & 98,6\% & -- \\
Fold 7 (S8)  & 352 & 35,8\% & 24,7\% & 91,4\% & 90,9\% & 0,0\% & 84,1\% & 97,8\% & 99,0\% & 0,0\% \\
Fold 8 (S9)  & 350 & 34,6\% & 8,6\% & 69,19\% & 64,8\% & 0,0\% & 87,4\% & 54,7\% & 84,3\% & 0,0\% \\
Fold 9 (S10) & 365 & 37,5\% & 16,2\% & 69,23\% & 98,6\% & 0,0\% & 14,8\% & 100,0\% & 100,0\% & 0,0\% \\
Fold 10 (S11)& 357 & 35,9\% & 32,2\% & 98,0\% & 100,0\% & 0,0\% & 94,4\% & 100,0\% & 100,0\% & 0,0\% \\
Fold 11 (S13)& 354 & 35,6\% & 42,1\% & 90,0\% & 89,0\% & 0,0\% & 63,0\% & 100,0\% & 100,0\% & 0,0\% \\
Fold 12 (S14)& 357 & 35,9\% & 20,2\% & 65,3\% & 24,7\% & -- & 64,1\% & 0,0\% & 96,1\% & -- \\
Fold 13 (S15)& 358 & 36,3\% & 30,4\% & 89,0\% & 90,8\% & 6,8\% & 44,7\% & 100,0\% & 100,0\% & 0,0\% \\
Fold 14 (S16)& 356 & 35,7\% & 11,5\% & 100,0\% & 100,0\% & 0,0\% & 95,5\% & 100,0\% & 100,0\% & 0,0\% \\
Fold 15 (S17)& 366 & 37,4\% & 28,1\% & 98,2\% & 89,3\% & 30,0\% & 76,5\% & 100,0\% & 100,0\% & \textbf{0,0\%} \\
\midrule
\textbf{Trung bình $\pm$ ĐLC} & \textbf{5294} & \textbf{35,5 $\pm$ 1,0\%} & \textbf{20,0 $\pm$ 10,5\%} & \textbf{85,1 $\pm$ 12,0\%} & \textbf{84,5 $\pm$ 19,9\%} & \textbf{14,3 $\pm$ 26,8\%} & \textbf{69,6 $\pm$ 22,2\%} & \textbf{83,1 $\pm$ 32,6\%} & \textbf{96,1 $\pm$ 7,1\%} & \textbf{6,6 $\pm$ 16,8\%} \\
\bottomrule
\end{tabular}}
\end{table*}

\begin{figure*}[t]
\centering
\includegraphics[width=0.98\linewidth]{Images/subject_loso_fold.png}
\caption{Phân phối hiệu năng theo từng đối tượng qua 15 fold LOSO tại $\tau = 0,80$: Điểm Selective Stress $F_{1,\text{sel}}$ (cột xanh) và Độ bao phủ hữu dụng (cột cam) trên từng cá nhân WESAD ($S2 \rightarrow S17$) cùng mức trung bình toàn bộ ($F_{1,\text{sel}} = 83,1\%$, Coverage $= 69,6\%$).}
\label{fig:subject_loso_fold}
\end{figure*}

Phân tích biến thiên cá nhân trong Bảng~\ref{tab:per_subject_loso} và Hình~\ref{fig:subject_loso_fold} cho thấy:
\begin{enumerate}
    \item Trên \textbf{8 trong số 15 folds} (S4, S6, S10, S11, S13, S15, S16, S17), phân loại chọn lọc đạt độ chính xác tuyệt đối $F_{1,\text{sel}} = 100,0\%$, và hai fold khác đạt kết quả rất cao (S7: 98,6\%, S8: 97,8\%).
    \item Độ bao phủ trung bình theo đối tượng đạt $69,6\% \pm 22,2\%$ (độ bao phủ gộp là $69,42\%$), dao động từ $14,8\%$ (S10) đến $100,0\%$ (S4), với độ chính xác chọn lọc trung bình $96,1\% \pm 7,1\%$. Sự sụt giảm độ bao phủ mạnh ở S10 ($14,8\%$) phản ánh sự thận trọng tối đa của mô hình trước độ bất định sinh lý cá nhân, sẵn sàng từ chối 85,2\% cửa sổ lấp lửng để bảo vệ độ chính xác chẩn đoán tuyệt đối ($F_{1,\text{sel}} = 100,0\%$, $\text{FAR}_{\text{sel}} = 0,0\%$). Điều này khẳng định hạn chế của một ngưỡng tĩnh toàn cục $\tau$ và gợi mở hướng nghiên cứu ngưỡng thích ứng theo cá nhân.
    \item Tỷ lệ báo động giả khi vận động $\text{FAR}_{\text{sel}}$ bị triệt tiêu về $0,0\%$ trên \textbf{11 trong số 13 đối tượng} có cửa sổ vận động, chỉ còn hai đối tượng còn báo động giả sót lại (S3: 54,5\%, S5: 31,6\%).
\end{enumerate}

\subsection{Nghiên cứu Cắt bỏ Đặc trưng Gia tốc (Trả lời RQ4 - Phần 1)}
Để kiểm chứng vai trò của cảm biến gia tốc ba trục, chúng tôi tiến hành ablation study qua 5 seed ngẫu nhiên (42, 123, 456, 789, 2024) trên ba cấu hình:
\begin{itemize}
    \item \textbf{Cấu hình A (Chỉ Sinh lý):} 19 đặc trưng (BVP, EDA, TEMP).
    \item \textbf{Cấu hình B (Đầy đủ Bối cảnh):} 25 đặc trưng (Sinh lý + 6 ACC).
    \item \textbf{Cấu hình C (Chỉ Gia tốc):} 6 đặc trưng gia tốc kế.
\end{itemize}

\begin{table}[htbp]
\caption{So sánh cắt bỏ đặc trưng (Ablation) qua 5 seed ngẫu nhiên tại $\tau = 0,80$ trên WESAD LOSO (Trung bình $\pm$ Độ lệch chuẩn).}
\label{tab:ablation_seeds}
\centering
\resizebox{\linewidth}{!}{
\begin{tabular}{lcccc}
\toprule
\textbf{Cấu hình} & \textbf{Độ Bao phủ (\%)} & \textbf{Sel. Macro $F_1$ (\%)} & $\text{FAR}_{\text{motion}}^{\text{sel}}$ (\%) & $\text{FAR}_{\text{motion}}^{\text{full}}$ (\%) \\
\midrule
A: Chỉ Sinh lý (Không ACC)  & $68,60 \pm 0,71$ & $92,98 \pm 0,61$ & $4,81 \pm 0,56$ & $13,69 \pm 0,74$ \\
\textbf{B: Đầy đủ (+ 6 ACC)} & $\mathbf{69,60 \pm 0,33}$ & $\mathbf{94,79 \pm 0,50}$ & $\mathbf{3,46 \pm 0,65}$ & $\mathbf{10,44 \pm 0,44}$ \\
C: Chỉ Gia tốc (ACC)        & $45,14 \pm 0,93$ & $82,38 \pm 0,66$ & $6,10 \pm 1,59$ & $26,31 \pm 0,70$ \\
\bottomrule
\end{tabular}}
\end{table}

\begin{figure}[htbp]
\centering
\includegraphics[width=\linewidth]{Images/ablation_acc_impact.png}
\caption{So sánh ablation đa chỉ số tại $\tau=0,80$: Bối cảnh gia tốc giúp gia tăng độ bao phủ đồng thời triệt tiêu báo động giả khi vận động.}
\label{fig:ablation_acc}
\end{figure}

Kết quả trong Bảng~\ref{tab:ablation_seeds} và Hình~\ref{fig:ablation_acc} chỉ ra:
\begin{enumerate}
    \item \textbf{Tính Ổn định Số học Qua các Seeds:} Bổ sung 6 đặc trưng ACC (Cấu hình B) cải thiện số học nhất quán so với chỉ dùng sinh lý (Cấu hình A) qua các seed. Điểm Selective Macro $F_1$ tăng từ $92,98\% \pm 0,61\%$ lên $94,79\% \pm 0,50\%$ ($+1,81\%$), trong khi $\text{FAR}_{\text{motion}}^{\text{sel}}$ giảm từ $4,81\% \pm 0,56\%$ xuống $3,46\% \pm 0,65\%$ (giảm $1,35\%$). Ở chế độ full coverage, ACC giúp giảm $\text{FAR}_{\text{motion}}^{\text{full}}$ từ $13,69\%$ xuống $10,44\%$ ($-3,25\%$).
    \item \textbf{Đánh giá Cấu hình Chỉ Dùng Gia tốc (Cấu hình C):} Cấu hình C (Chỉ dùng ACC) đạt Selective Macro $F_1 = 82,38\% \pm 0,66\%$ (Stress $F_1 = 74,23\% \pm 1,11\%$), khớp hoàn toàn với Hình~\ref{fig:ablation_acc}. Điều này giải thích bởi giao thức WESAD: pha TSST yêu cầu đứng thuyết trình và làm toán nhẩm nên người tham gia cử động tay nhiều hơn hẳn so với lúc ngồi yên đọc sách hoặc thiền. Tuy nhiên, Cấu hình C bị sụt giảm độ bao phủ nghiêm trọng ($45,14\%$) và có tỷ lệ báo động giả toàn phần rất cao ($\text{FAR}_{\text{motion}}^{\text{full}} = 26,31\%$), khẳng định gia tốc kế không thể thay thế tín hiệu sinh lý tự chủ. Điều này giải thích bởi giao thức WESAD: pha TSST yêu cầu đứng thuyết trình và làm toán nhẩm nên người tham gia cử động tay nhiều hơn hẳn so với lúc ngồi yên đọc sách (baseline) hoặc thiền. Do đó, gia tốc mang tín hiệu tương quan với kịch bản phòng lab. Tuy nhiên, Cấu hình C có tỷ lệ báo động giả cao ($\text{FAR}_{\text{motion}}^{\text{sel}} = 8,47\%$), khẳng định gia tốc kế không thể thay thế tín hiệu sinh lý mà chỉ đóng vai trò cung cấp bối cảnh động học bổ trợ.
    \item \textbf{Chi tiết Kiểm định Thống kê và Số Cặp Hiệu dụng:} Kiểm định Wilcoxon signed-rank trên 15 fold (seed 42) ghi nhận: $F_{1,\text{sel}}$ ($W = 9,0, p = 0,753, r_{rb} = 0,143$), $\text{FAR}_{\text{sel}}$ ($W = 1,0, p = 1,000, r_{rb} = -0,333$), và Coverage ($W = 37,0, p = 0,552, r_{rb} = 0,187$), với $p_{\text{holm}} = 1,000$. Khoảng tin cậy bootstrap cho chênh lệch $(B - A)$ chứa số không: $\Delta F_{1,\text{sel}} = +0,0342$ (95\% CI: $[-0,0278; +0,1469]$) và $\Delta \text{FAR}_{\text{sel}} = -1,0169\%$ (95\% CI: $[-4,5977; +1,0870]$). Lý do $p > 0,05$ là vì trên nhiều đối tượng, cả hai cấu hình đều đạt $0\%$ báo động giả (chỉ có $n_{\text{nonzero}} = 2$ cặp khác biệt cho $\text{FAR}_{\text{sel}}$ và $n_{\text{nonzero}} = 6$ cho $F_1$), khiến kiểm định Wilcoxon bị giới hạn về lực thống kê trên mẫu $N=15$.
\end{enumerate}

\subsection{Đóng góp Đặc trưng TreeSHAP và Phân tích Ngoại lai (Trả lời RQ4 - Phần 2)}

\textbf{1) Đóng góp Tổng thể theo Cảm biến:} Phân tích TreeSHAP \cite{lundberg2020local} trên toàn bộ dự đoán ngoài tập (Hình~\ref{fig:shap}) chỉ ra EDA chiếm $31,6\%$ đóng góp, tiếp theo là BVP/HRV ($26,6\%$), TEMP ($25,9\%$), và ACC ($15,9\%$). Bốn đặc trưng quan trọng nhất là \texttt{eda\_tonic\_mean} (13,6\%), \texttt{hrv\_mean\_hr} (12,3\%), \texttt{temp\_max} (10,6\%), và \texttt{eda\_scr\_peaks\_count} (10,3\%).

\begin{figure}[htbp]
\centering
\includegraphics[width=0.98\linewidth]{Images/shap_summary.png}
\caption{Đóng góp thuộc tính đặc trưng toàn cục TreeSHAP cho mô hình nhận diện stress từ cổ tay.}
\label{fig:shap}
\end{figure}

\textbf{2) Mẫu hình Xung đột Tín hiệu Kích hoạt Từ chối:} Phân tích các cửa sổ bị từ chối (\textit{Bất định}) cho thấy: khi gia tốc ghi nhận rung lắc mạnh trong khi nhịp tim và dẫn truyền da dao động không dứt khoát, sự xung đột này khiến xác suất hậu nghiệm rơi vào vùng lưỡng lự $[1-\tau, \tau]$ và kích hoạt từ chối phán đoán an toàn.

\textbf{3) Phân ly Sinh lý ở Ca Ngoại lai Subject S14:} Hình~\ref{fig:failure_s12} phân tích trường hợp cá nhân S14 (Fold 12). Dữ liệu cho thấy S14 thuộc nhóm \textit{giảm phản ứng điện da} (electrodermal hypo-responder) \cite{boucsein2012electrodermal}: Điểm z-score của \texttt{eda\_tonic\_mean} ($z = -0,90$) và \texttt{eda\_scr\_peaks\_count} ($z = -1,02$) rất thấp, khiến riêng đặc trưng \texttt{eda\_tonic\_mean} tạo ra đóng góp SHAP âm mạnh ($-3,22$) kéo dự đoán về Không Stress, áp đảo hoàn toàn tín hiệu nhịp tim nhanh ($z = +2,76$). Sự phân ly sinh lý này giải thích vì sao mô hình cây dựa nhiều vào EDA bị thất bại ($F_{1,\text{sel}} = 0,0\%$) trên người có cơ địa tiết mồ hôi không điển hình, trong khi mạng MLP lại khai thác được các kết hợp phi tuyến ($F_1 = 97,2\%$). Hiện tượng này nhấn mạnh sự cần thiết của cơ chế thích ứng cá nhân hóa trong các nghiên cứu tiếp theo.

\begin{figure}[htbp]
\centering
\includegraphics[width=0.98\linewidth]{Images/fig_failure_S12.png}
\caption{Hiện tượng phân ly sinh lý trên ca ngoại lai S14: Phản ứng điện da kém nhạy lấn át tín hiệu nhịp tim nhanh trong pha căng thẳng.}
\label{fig:failure_s12}
\end{figure}

% =========================================================================
% VI. HẠN CHẾ VÀ CÂN NHẮC THỰC TIỄN
% =========================================================================
\section{Hạn chế và Cân nhắc Thực tiễn}
Mặc dù khung phương pháp đạt được hiệu quả vượt trội trong việc triệt tiêu báo động giả, các hạn chế phương pháp luận sau đây cần được nhìn nhận minh bạch:
\begin{enumerate}
    \item \textbf{Quy mô Một Bộ Dữ liệu Phòng Thí nghiệm ($N=15$):} WESAD là tập dữ liệu thu thập trong môi trường phòng lab có kiểm soát với 15 người dưới kích thích giao thức TSST. Căng thẳng tâm lý trong đời sống thực phong phú và phức tạp hơn. Việc kiểm chứng mô hình trên các tập dữ liệu tự do ngoài đời thực là bước đi cần thiết tiếp theo.
    \item \textbf{Mối Tương quan Gắn liền Giao thức Thí nghiệm:} Trong WESAD, pha TSST đòi hỏi người tham gia đứng thuyết trình và làm tính nhẩm trước ban giám khảo, tự nhiên tạo ra nhiều cử động tay hơn so với lúc ngồi đọc sách hoặc thiền. Do đó, gia tốc mang tính tương quan với kịch bản lab. Việc khái quát hóa sang các hoạt động thể thao mạnh (chạy bộ, tập gym) đòi hỏi các tập dữ liệu có bài tập thể chất riêng biệt.
    \item \textbf{Ngưỡng Tin cậy Tĩnh Toàn cục ($\tau = 0,80$):} Mặc dù $\tau = 0,80$ đem lại điểm cân bằng tối ưu trên dữ liệu gộp, độ bao phủ thực tế trên từng người dao động từ 14,8\% (S10) đến 100,0\% (S4). Độ bao phủ thấp cục bộ cho thấy một ngưỡng tĩnh có thể quá khắt khe với một số cá nhân, gợi mở hướng phát triển ngưỡng thích ứng động theo cơ địa sinh lý của từng người.
    \item \textbf{Độ Phức tạp Tính toán Trích xuất Đặc trưng Nhúng:} Mặc dù suy luận của cây quyết định rất nhẹ ($< 1$\,ms), khâu trích xuất đặc trưng sinh lý—đặc biệt là bài toán tối ưu lồi trong cvxEDA để tách sóng EDA và biến đổi FFT cho HRV—đòi hỏi tài nguyên tính toán đáng kể trên vi điều khiển năng lượng siêu thấp. Các thuật toán xấp xỉ số nguyên nhẹ là hướng tối ưu hóa phần cứng quan trọng trong tương lai.
\end{enumerate}

% =========================================================================
% VII. KẾT LUẬN
% =========================================================================
\section{Kết luận}
Nghiên cứu này đã so sánh một cách hệ thống giữa bài toán nhận diện căng thẳng bắt buộc 100\% và cơ chế phân loại chọn lọc có từ chối trên thiết bị đeo cổ tay. Qua khảo sát tám kiến trúc học máy dưới giao thức 15-fold LOSO trên WESAD, hiện tượng báo động giả do vận động chứng minh là lỗ hổng mang tính hệ thống dưới chế độ ép buộc phán đoán ($\text{FAR}_{\text{motion}} \approx 12,2\% - 24,0\%$), trong khi phân loại chọn lọc triệt tiêu từ 55,0\% đến 90,0\% các cảnh báo sai này. Hệ thống đề xuất đạt các chỉ số độ tin cậy y tế xuất sắc (Selective Macro $F_1 = 95,20\%$, ECE $= 0,0290$ trên tập giữ lại, $\text{FAR}_{\text{motion}} = 4,07\%$) tại độ bao phủ vận hành $69,42\%$. Được củng cố bởi khả năng giải thích TreeSHAP và kiểm chứng trên 15 đối tượng, công trình đặt ra các nguyên lý thiết kế thực tiễn cho các hệ thống giám sát sức khỏe đeo tay đáng tin cậy.

% =========================================================================
% TÀI LIỆU THAM KHẢO (38 TÀI LIỆU CHUẨN IEEE)
% =========================================================================
\begin{thebibliography}{00}
\bibitem{cohen2007psychological}
S.~Cohen, D.~Janicki-Deverts, and G.~E.~Miller, ``Psychological stress and disease,'' \textit{JAMA}, vol.~298, no.~14, pp. 1685--1687, 2007.

\bibitem{schneiderman2005stress}
N.~Schneiderman, G.~Ironson, and S.~D.~Siegel, ``Stress and health: Psychological, behavioral, and biological determinants,'' \textit{Annu. Rev. Clin. Psychol.}, vol.~1, pp. 607--628, 2005.

\bibitem{can2019stress}
Y.~S.~Can, B.~Arnrich, and C.~Ersoy, ``Stress detection in daily life scenarios using smart phones and wearable devices: A review,'' \textit{J. Biomed. Inform.}, vol.~92, p. 103139, 2019.

\bibitem{gjoreski2016continuous}
M.~Gjoreski, H.~Gjoreski, M.~Lu\v{s}trek, and M.~Gams, ``Continuous stress detection using a wrist device: In laboratory and real life,'' in \textit{Proc. ACM Int. Joint Conf. Pervasive Ubiquitous Comput. (UbiComp Adjunct)}, 2016, pp. 1185--1193.

\bibitem{schmidt2018wesad}
P.~Schmidt, A.~Reiss, R.~Duerichen, C.~Marberger, and K.~Van~Laerhoven, ``Introducing WESAD, a multimodal dataset for wearable stress and affect detection,'' in \textit{Proc. 20th ACM Int. Conf. Multimodal Interaction (ICMI)}, 2018, pp. 400--408.

\bibitem{sano2018multimodal}
A.~Sano, A.~J.~Phillips, A.~Z.~Yu, A.~W.~McHill, S.~Taylor, N.~Jaques, C.~A.~Czeisler, E.~B.~Klerman, and R.~W.~Picard, ``Recognizing academic performance, sleep and stress from objective mobile, wearable, and network data,'' \textit{IEEE Trans. Affect. Comput.}, vol.~14, no.~2, pp. 977--990, 2023.

\bibitem{schumm2008effect}
J.~Schumm, M.~B\"achlin, C.~Setz, B.~Arnrich, D.~Roggen, and G.~Tr\"oster, ``Effect of movement on the electrodermal response after a startle event,'' \textit{Methods Inf. Med.}, vol.~47, no.~3, pp. 186--191, 2008.

\bibitem{tamura2014wearable}
T.~Tamura, Y.~Maeda, M.~Sekine, and M.~Yoshida, ``Wearable photoplethysmographic sensors---past and present,'' \textit{Electronics}, vol.~3, no.~2, pp. 282--302, 2014.

\bibitem{zhang2014troika}
Z.~Zhang, Z.~Pi, and B.~Liu, ``TROIKA: A general framework for heart rate monitoring using wrist-type photoplethysmographic signals during intensive physical exercise,'' \textit{IEEE Trans. Biomed. Eng.}, vol.~62, no.~2, pp. 522--531, 2015.

\bibitem{boucsein2012electrodermal}
W.~Boucsein, \textit{Electrodermal Activity}, 2nd~ed. New York, NY, USA: Springer Science \& Business Media, 2012.

\bibitem{siirtola2019continuous}
P.~Siirtola, ``Continuous stress detection using the sensors of commercial smartwatch,'' in \textit{Proc. ACM Int. Joint Conf. Pervasive Ubiquitous Comput. (UbiComp/ISWC Adjunct)}, 2019, pp. 1198--1201.

\bibitem{bobade2020stress}
P.~Bobade and M.~Vani, ``Stress detection with machine learning and deep learning using multimodal physiological data,'' in \textit{Proc. 2nd Int. Conf. Inventive Research in Computing Applications (ICIRCA)}, 2020, pp. 51--57.

\bibitem{dalmeida2021hrv}
K.~M.~Dalmeida and G.~L.~Masala, ``HRV features as viable physiological markers for stress detection using wearable devices,'' \textit{Sensors}, vol.~21, no.~8, p. 2873, 2021.

\bibitem{behinaein2021transformer}
B.~Behinaein, A.~Bhatti, D.~Rodenburg, P.~Hungler, and A.~Etemad, ``A transformer architecture for stress detection from ECG,'' in \textit{Proc. ACM Int. Symp. Wearable Comput. (ISWC)}, 2021, pp. 132--134.

\bibitem{bottcher2023detection}
K.~B\"ottcher, S.~Poser, F.~Heinrich, and K.~Van~Laerhoven, ``Detection of mental stress using consumer smartwatches,'' \textit{IEEE Trans. Affect. Comput.}, vol.~14, no.~4, pp. 2704--2716, 2023.

\bibitem{guo2017calibration}
C.~Guo, G.~Pleiss, Y.~Sun, and K.~Q.~Weinberger, ``On calibration of modern neural networks,'' in \textit{Proc. 34th Int. Conf. Machine Learning (ICML)}, 2017, pp. 1321--1330.

\bibitem{sendelbach2013alarm}
S.~Sendelbach and M.~Funk, ``Alarm fatigue: A patient safety concern,'' \textit{AACN Adv. Crit. Care}, vol.~24, no.~4, pp. 378--386, 2013.

\bibitem{chow1970optimum}
C.~K.~Chow, ``On optimum recognition error and reject tradeoff,'' \textit{IEEE Trans. Inf. Theory}, vol.~16, no.~1, pp. 41--46, 1970.

\bibitem{cortes2016learning}
C.~Cortes, G.~DeSalvo, and M.~Mohri, ``Learning with rejection,'' in \textit{Proc. Int. Conf. Algorithmic Learning Theory (ALT)}, 2016, pp. 67--82.

\bibitem{geifman2017selective}
Y.~Geifman and R.~El-Yaniv, ``Selective classification for deep neural networks,'' in \textit{Advances in Neural Information Processing Systems (NeurIPS)}, vol.~30, 2017, pp. 4878--4887.

\bibitem{geifman2019selectivenet}
Y.~Geifman and R.~El-Yaniv, ``SelectiveNet: A deep neural network with an integrated reject option,'' in \textit{Proc. 36th Int. Conf. Machine Learning (ICML)}, 2019, pp. 2151--2159.

\bibitem{grandvalet2008support}
Y.~Grandvalet, A.~Rakotomamonjy, J.~Keshet, and S.~Canu, ``Support vector machines with a reject option,'' in \textit{Advances in Neural Information Processing Systems (NeurIPS)}, vol.~21, 2008, pp. 537--544.

\bibitem{hendrickx2024machine}
K.~Hendrickx, L.~Perini, D.~Van~der~Plas, W.~Meert, and J.~Davis, ``Machine learning with a reject option: A survey,'' \textit{Mach. Learn.}, vol.~113, no.~5, pp. 3073--3110, 2024.

\bibitem{saeb2017need}
S.~Saeb, L.~Lonini, A.~Jayaraman, D.~C.~Mohr, and K.~P.~Kording, ``The need to approximate the use-case in clinical machine learning,'' \textit{GigaScience}, vol.~6, no.~5, p. gix019, 2017.

\bibitem{little2017using}
M.~A.~Little, G.~Varoquaux, S.~Saeb, L.~Lonini, A.~Jayaraman, D.~C.~Mohr, and K.~P.~Kording, ``Using and understanding cross-validation strategies: Perspectives on Saeb et al.,'' \textit{GigaScience}, vol.~6, no.~5, p. gix020, 2017.

\bibitem{lundberg2020local}
S.~M.~Lundberg, G.~Erion, H.~Chen, A.~DeGrave, J.~M.~Prutkin, B.~Nair, R.~Katz, J.~Himmelfarb, N.~Bansal, and S.-I.~Lee, ``From local explanations to global understanding with explainable AI for trees,'' \textit{Nature Machine Intelligence}, vol.~2, no.~1, pp. 56--67, 2020.

\bibitem{taskforce1996heart}
Task Force of the European Society of Cardiology and the North American Society of Pacing and Electrophysiology, ``Heart rate variability: Standards of measurement, physiological interpretation, and clinical use,'' \textit{Circulation}, vol.~93, no.~5, pp. 1043--1065, 1996.

\bibitem{shaffer2017overview}
F.~Shaffer and J.~P.~Ginsberg, ``An overview of heart rate variability metrics and norms,'' \textit{Front. Public Health}, vol.~5, p. 258, 2017.

\bibitem{garbarino2014empatica}
M.~Garbarino, M.~Lai, D.~Bender, R.~W.~Picard, and S.~Tognetti, ``Empatica E3---A wearable wireless multi-sensor device for real-time computerized biofeedback and data acquisition,'' in \textit{Proc. 4th Int. Conf. Wirel. Mobile Commun. Healthc. (MobiHealth)}, 2014, pp. 39--42.

\bibitem{platt1999probabilistic}
J.~Platt, ``Probabilistic outputs for support vector machines and comparisons to regularized likelihood methods,'' \textit{Adv. Large Margin Classif.}, vol.~10, no.~3, pp. 61--74, 1999.

\bibitem{kull2017beta}
M.~Kull, T.~Silva~Filho, and P.~Flach, ``Beta calibration: a well-founded and easily implemented alternative to logistic calibration for binary classifiers,'' in \textit{Proc. 20th Int. Conf. Artif. Intell. Stat. (AISTATS)}, 2017, pp. 623--631.

\bibitem{zadrozny2002transforming}
B.~Zadrozny and C.~Elkan, ``Transforming classifier scores into accurate multiclass probability estimates,'' in \textit{Proc. 8th ACM SIGKDD Int. Conf. Knowl. Discov. Data Min. (KDD)}, 2002, pp. 694--699.

\bibitem{niculescu2005predicting}
A.~Niculescu-Mizil and R.~Caruana, ``Predicting good probabilities with supervised learning,'' in \textit{Proc. 22nd Int. Conf. Machine Learning (ICML)}, 2005, pp. 625--632.

\bibitem{chen2016xgboost}
T.~Chen and C.~Guestrin, ``XGBoost: A scalable tree boosting system,'' in \textit{Proc. 22nd ACM SIGKDD Int. Conf. Knowl. Discov. Data Min. (KDD)}, 2016, pp. 785--794.

\bibitem{shwartzziv2022tabular}
R.~Shwartz-Ziv and R.~Armon, ``Tabular data: Deep learning is not all you need,'' \textit{Inf. Fusion}, vol.~81, pp. 84--90, 2022.

\bibitem{grinsztajn2022why}
L.~Grinsztajn, E.~Oyallon, and G.~Varoquaux, ``Why do tree-based models still outperform deep learning on typical tabular data?,'' in \textit{Advances in Neural Information Processing Systems (NeurIPS)}, vol.~35, 2022, pp. 507--520.

\bibitem{borisov2024deep}
V.~Borisov, T.~Leemann, K.~Se\ss{}ler, J.~Haug, M.~Pawelczyk, and G.~Kasneci, ``Deep neural networks and tabular data: A survey,'' \textit{IEEE Trans. Neural Netw. Learn. Syst.}, vol.~35, no.~6, pp. 7499--7519, 2024.

\bibitem{greco2016cvxeda}
A.~Greco, G.~Valenza, A.~Lanata, E.~P.~Scilingo, and R.~Citi, ``cvxEDA: A convex optimization approach to electrodermal activity processing,'' \textit{IEEE Trans. Biomed. Eng.}, vol.~63, no.~4, pp. 797--804, 2016.

\bibitem{tarvainen2014kubios}
M.~P.~Tarvainen, J.-P.~Niskanen, J.~A.~Lipponen, P.~O.~Ranta-Aho, and P.~A.~Karjalainen, ``Kubios HRV--heart rate variability analysis software,'' \textit{Comput. Methods Programs Biomed.}, vol.~113, no.~1, pp. 210--220, 2014.
\end{thebibliography}

\end{document}
